\section{Optimasi Performa Aplikasi Web: Caching, Minifikasi, Lazy Loading}

Optimasi performa meningkatkan kecepatan dan responsivitas aplikasi. Tiga teknik utama: \textbf{caching} menyimpan data yang sering diakses agar tidak perlu diproses ulang (browser cache, server cache, opcode cache); \textbf{minifikasi} mengurangi ukuran file CSS, JavaScript, dan HTML dengan menghapus whitespace, komentar, dan memperpendek nama variabel; \textbf{lazy loading} menunda pemuatan resource (gambar, video) yang belum terlihat di viewport, mengurangi waktu loading awal \cite{mdnToolsTesting}.

\begin{webcode}[language=JavaScript, caption={Lazy loading gambar dengan atribut HTML dan Intersection Observer}]
<!-- HTML: lazy loading native -->
<!-- <img src="hero.jpg" loading="lazy" alt="Hero image"> -->

// JavaScript: Intersection Observer untuk lazy loading kustom
const images = document.querySelectorAll('img[data-src]');
const observer = new IntersectionObserver((entries) => {
  entries.forEach(entry => {
    if (entry.isIntersecting) {
      const img = entry.target;
      img.src = img.dataset.src;
      img.removeAttribute('data-src');
      observer.unobserve(img);
    }
  });
});
images.forEach(img => observer.observe(img));
\end{webcode}

\begin{catatan}
Checklist optimasi performa: (1) aktifkan \textbf{gzip/brotli compression} di web server; (2) minifikasi CSS dan JS; (3) gunakan \code{loading="lazy"} pada gambar; (4) aktifkan browser caching dengan header \code{Cache-Control}; (5) optimalkan gambar (format WebP, kompresi); (6) gunakan CDN untuk aset statis; (7) aktifkan \textbf{opcode cache} (OPcache) untuk PHP.
\end{catatan}

